% 2021 全国甲卷文科 7：正视图与四个候选侧视图
% 绘制依据：源 PDF 第 1 页原图。五个同样大小的正方形。
% 正视图：一条【虚线】自左上区域降至右下角（斜率约 −1/2）。
% 四个选项的判别特征为「线型 × 走向」：
%   (A) 实线上升（自底边中部升至右边）  (B) 虚线下降（自左边降至底边右部）
%   (C) 虚线上升（自底边中部升至右边）  (D) 实线下降（自左边降至底边右部）
% 选项标签 (A)~(D) 置于各正方形左下角外侧。
\begin{tikzpicture}[font=\small, line width=0.4pt]
  \def\s{1.75}% 正方形边长
  % ---------- 正视图 ----------
  \begin{scope}[shift={({1.6*\s},{1.75*\s})}]
    \draw (0,0) rectangle (\s,\s);
    \draw[dashed] ({0.20*\s},{0.40*\s}) -- ({0.98*\s},{0.02*\s});
    \node[anchor=north] at ({0.5*\s},-0.10) {正视图};
  \end{scope}
  % ---------- 四个选项 ----------
  % (A) 实线上升
  \begin{scope}[shift={(0,0)}]
    \draw (0,0) rectangle (\s,\s);
    \draw ({0.42*\s},0) -- (\s,{0.60*\s});
    \node[anchor=east] at (-0.06,{0.06*\s}) {(A)};
  \end{scope}
  % (B) 虚线下降
  \begin{scope}[shift={({1.55*\s},0)}]
    \draw (0,0) rectangle (\s,\s);
    \draw[dashed] (0,{0.62*\s}) -- ({0.88*\s},0);
    \node[anchor=east] at (-0.06,{0.06*\s}) {(B)};
  \end{scope}
  % (C) 虚线上升
  \begin{scope}[shift={({3.10*\s},0)}]
    \draw (0,0) rectangle (\s,\s);
    \draw[dashed] ({0.42*\s},0) -- (\s,{0.60*\s});
    \node[anchor=east] at (-0.06,{0.06*\s}) {(C)};
  \end{scope}
  % (D) 实线下降
  \begin{scope}[shift={({4.65*\s},0)}]
    \draw (0,0) rectangle (\s,\s);
    \draw (0,{0.62*\s}) -- ({0.88*\s},0);
    \node[anchor=east] at (-0.06,{0.06*\s}) {(D)};
  \end{scope}
\end{tikzpicture}
